\chapter{Lambda-BLL}
\lhead[\fancyplain{}{\bfseries\thepage}]{\fancyplain{}{\bfseries\rightmark}}

This chapter introduces the $\lambda$-BLL calculus~\cite{lago2024computationalindistinguishabilitylogicalrelations}, a typed probabilistic $\lambda$-calculus to model polynomial-time cryptographic computation. Its design principles and language features capture the essential for expressing cryptographic experiments, adversaries, and reductions.

$\lamBLL$ provides a programming-language framework in which computational complexity is enforced statically by the type system. In particular, it supports probabilistic polynomial-time computation, bounded interaction with oracles, shared mutable state and higher-order reasoning about cryptographic constructions. These features make it well-suited for formalizing experiment-based cryptographic proofs.

The type system of $\lambda$-BLL is linear and graded, with a close correspondence to Bounded Linear Logic and graded calculi. Resource usage is tracked by polynomial grades, generated by:
\begin{align*}
p ::= 1 \mid i \mid p + p \mid p * p
\end{align*}
where the variable $i$ represents the security parameter.

The calculus follows a call-by-push-value (CBPV) discipline~\cite{levy2012call}, distinguishing between positive types and general types. This separation permits to restrict function application to positive arguments only.

The base of the type system consists of ground types:
\begin{align*}
\UU \mid \BB \mid \SS[p] 
\end{align*}
representing unit, booleans, and fixed-length bitstrings of polynomially bounded size. These types correspond to the fundamental data manipulated in cryptographic protocols, such as nonces, messages, and keys.

Positive types extend ground types with:
\begin{align*}
P \otimes P \mid \bang_p A 
\end{align*}
namely pairs formed with the tensor products $P \otimes P$, and the graded comonadic modality $\bang_p A$. In a linear type system, a value must normally be used exactly once. The modality $\bang_p A$ relaxes linearity by allowing a value of type $A$ to be used up to $p$ times. This is central to the cryptographic interpretation of the calculus because it enables a static enforcement of polynomial-time bounds. For example, an oracle of type $\bang_{q(i)} \Oracle$ may be queried at most $q(i)$ times by an adversary for a polynomial $q$. Any attempt to exceed this bound results in an ill-typed term, ensuring that polynomial-time constraints are enforced by typing rather than by operational reasoning. As a result, the proof fails at the type level rather than at runtime.

\begin{figure}[htpb]
  \centering
  \begin{align*}
    G &::= \UU \mid \BB \mid \SS[p] \tag{Ground types} \\
    P &::= G \mid P \otimes P \mid \bang_p A \tag{Positive types} \\
    A &::= P \mid P \multimap A \tag{Types}
  \end{align*}
  \caption{Types of $\lambda$-BLL}
  \label{fig:types}
\end{figure}

$\lamBLL$ extends the pure lambda-calculus with features necessary for modelling cryptographic computation. The language define a collection of function symbols $\mathcal{F}$. Each function symbol $f \in \mathcal{F}$ has a type denoted by $\mathtt{typeof}$ of the form
\begin{align*}
\mathtt{typeof}(f): G_1 \times \dots \times G_m \to G
\end{align*}
where $G_1, \dots, G_m$ and $G$ are ground types. For each polynomial $p$, a corresponding constructor $f_p$ represents a resource-bounded invocation of $f$, with the grade $p$ accounting for its computational cost within the type system.

Values represent data that can be passed, stored or paired but not directly executed, while computations represent effectful execution. Ground values include $\star$, booleans, and bitstrings. Positive values include variables, pairs, and suspended computations of the form $!M$. Computations follow the CBPV structure so values are injected using $\return$, suspended computations are forced using $\der$, and applications is defined only on positive arguments.

\begin{figure}[htpb]
  \centering
  \begin{align*}
    W &::= \star \mid \true \mid \false \mid s \tag{Ground values} \\
    Z &::= x \mid W\mid \lis{Z,Z} \mid \bang M \tag{Positive values} \\
    \Val \ni V &::= Z \mid \lambda x.M \tag{Values} \\
    \Lambda \ni M &::=  \return V \mid  \der (Z) \mid  MZ \tag{Computations} \\
    &\mid \letin{x=N}{M} | \letin{\lis{x,y}=Z}{M} \\
    &\mid \ifthen{Z}{M}{N} \mid \loopp{V}{p}{M} \\
    &\mid \setr{r}{Z} \mid \get r \\
    &\mid f_p(Z_1, \dots, Z,m) \\
  \end{align*}
  \caption{Syntax of $\lamBLL$}
  \label{fig:syntax}
\end{figure}

To model cryptographic experiments involving shared state, $\lamBLL$ includes global references. References store ground values (and ground only) and are manipulated via $\setr{r}{Z}$ and $\get r$. This ensures that stateful effects do not introduce higher-order behavior beyond what is explicitly admitted by the type system. The language provides also control constructs such as conditionals, let-bindings, and pattern matching on pairs. Bounded iteration is expressed using the construct $\loopp{V}{p}{M}$, which executes the computation $M$ at most $p$ times.

Since variables in $\lamBLL$ are classified according to their types, a variable is said to be ground if its type contains no function arrows, preventing ground terms from encoding higher-order behavior thus being dangerous as global references. To track higher-order dependencies, the calculus distinguishes the set of free higher-order variables of a term $M$, denoted $\mathit{HOFV}(M)$. This set contains exactly those free variables whose types include at least one occurrence of the function type constructor $(\multimap)$.

$\lamBLL$ considers two different typing contexts
\begin{align*}
\Gamma &::= \emptyset \mid \Gamma, x: P \tag{Variable context} \\
\Theta &::= \emptyset \mid \Theta, r: P \tag{Reference context}
\end{align*}
The context $\Gamma$ contains term variables with positive types only, which are not computations, while the reference context $\Theta$ contains references which are can only be ground types. Fixed a reference context $\Theta$ then there are two typing judgments with respect to $\Theta$
\begin{align*}
\Gamma &\vdash_v^\Theta V: A \tag{Value typing} \\
\Gamma &\vdash_c^\Theta M: A \tag{Computation typing}
\end{align*}
Computations may access references, but references never contain computations.
% % \begin{figure}[h]
% % 	\centering
% % 	\begin{tikzpicture}[]
% % 		\node  at (0,0) {$\mathtt{false}$};
% % 		\node[rotate=45] (A) at (-1,-0.5) {$\not\coarseq$};
% % 		\node[rotate=-45]  at (1,-0.5) {$\coarseq$};
% % 		\node  at (-2, -1) {$\mathtt{MacForge}^F$};
% % 		\node  at (2, -1) {$\widehat{\mathtt{MacForge}}$};
% % 		\node[rotate=90]  at (-2,-1.5) {$\approx$};
% % 		\node[rotate=90]  at (2,-1.5) {$\approx$};
% % 		\node  at (-2, -2) {$D^F$};
% % 		\node  at (2, -2) {$\widehat{D}$};
% % 		\node (B) at (0,-2) {$\not\coarseq$};
% % 		\draw[-{Implies},double] (A) to (B);
% % 	\end{tikzpicture}
% % 	\caption{Outline Proof of Security for $\mathtt{MacForge}$}\label{fig:proofscheme}
% % \end{figure}
